\section*{Chapter \ref{ch:g10:measuring-the-earth} --- Measuring the Earth}

\begin{solution}{exo:g10:measuring-the-earth:1}
A ship disappears hull first: over a flat sea it would only shrink, all of
it remaining visible. The Earth's shadow on the Moon in an eclipse is always
an arc of a circle: a flat disc would cast an elongated shadow when seen
edge-on. Travelling north, the Pole Star rises above the horizon: on a flat
Earth its height would stay almost the same.
\end{solution}

\begin{solution}{exo:g10:measuring-the-earth:2}
\omterm{def:g10:measuring-the-earth:geodesy}{Geodesy} measures the shape and size of the Earth and the positions of
points on it, and their changes with time. \omterm{def:g10:measuring-the-earth:triangulation}{Triangulation} locates points by
a chain of triangles in which all the angles are measured; only one length,
the baseline, is measured on the ground.
\end{solution}

\begin{solution}{exo:g10:measuring-the-earth:3}
$C = 50 \times 5000 = 250\,000$ stadia; $250\,000 \times \qty{157.2}{m} =
\qty{39300}{km}$.
\end{solution}

\begin{solution}{exo:g10:measuring-the-earth:4}
$111/60 \approx \qty{1.85}{km}$. Quarter meridian: $90 \times 111 \approx
\qty{10000}{km}$.
\end{solution}

\begin{solution}{exo:g10:measuring-the-earth:5}
$a - b = 6378.137 - 6356.752 = \qty{21.385}{km}$;
$f = 21.385/6378.137 \approx 0.003353 \approx 1/298$.
\end{solution}

\begin{solution}{exo:g10:measuring-the-earth:6}
South: $\tan\theta_1 = 0.364$, $\theta_1 \approx \qty{20.0}{\degree}$. North:
$\tan\theta_2 = 0.456$, $\theta_2 \approx \qty{24.5}{\degree}$.
$\alpha = \qty{4.5}{\degree}$, so $C = (360/4.5) \times 500 = 80 \times 500 =
\qty{40000}{km}$.
\end{solution}

\begin{solution}{exo:g10:measuring-the-earth:7}
$C = 180 - 62 - 71 = \qty{47}{\degree}$.
$AC = 10.00 \sin 71^\circ/\sin 47^\circ = 10.00 \times 0.9455/0.7314 \approx
\qty{12.93}{km}$; $BC = 10.00 \sin 62^\circ/\sin 47^\circ \approx
\qty{12.07}{km}$.
\end{solution}

\begin{solution}{exo:g10:measuring-the-earth:8}
The change of longitude is $10 + 20 = \qty{30}{\degree}$. At
\qty{60}{\degree} N a degree is $111.3 \times 0.5 = \qty{55.65}{km}$, so the
distance is $30 \times 55.65 \approx \qty{1670}{km}$. On the equator:
$30 \times 111.3 \approx \qty{3340}{km}$, twice as far.
\end{solution}

\begin{solution}{exo:g10:measuring-the-earth:9}
The line of sight touches the sea at $T$: a tangent to a circle is
perpendicular to the radius at the point of contact, so the angle at $T$ is
a right angle. Pythagoras: $(R+h)^2 = R^2 + d^2$, so $d^2 = 2Rh + h^2$; $h$ is
tiny compared with $R$, so $d \approx \sqrt{2Rh} = \sqrt{2 \times
6\,371\,000 \times 2} \approx \qty{5050}{m}$, about \qty{5}{km}.
\end{solution}

\begin{solution}{exo:g10:measuring-the-earth:10}
The curve is steepest in the middle latitudes, around \qty{30}{\degree} to
\qty{60}{\degree}, and flat near the equator and the pole. On a flattened
Earth the surface is less curved near the poles; the less curved the surface,
the farther one must travel for the vertical (and so the latitude) to turn by
one degree.
\end{solution}

\begin{solution}{exo:g10:measuring-the-earth:11}
$1072.2/9.646 \approx \qty{111.15}{km}$ per degree; $90 \times 111.15
\approx \qty{10004}{km}$, \qty{2}{km} more than the true quarter meridian.
The arc is centred near \qty{46}{\degree}, where a degree is slightly longer
than the average degree from the equator to the pole
($10\,002.0/90 \approx \qty{111.13}{km}$).
\end{solution}

\begin{solution}{exo:g10:measuring-the-earth:12}
With \qty{7.1}{\degree}: $C = (360/7.1) \times 5000 \approx 253\,500$ stadia
instead of $250\,000$: about 1.4\,\% too large (a smaller angle for the same
distance means a bigger Earth). The distance, paced by walkers along the
Nile, was much less precise than the angle.
\end{solution}

\begin{solution}{exo:g10:measuring-the-earth:13}
In $ABC$: $C = \qty{50}{\degree}$, $AC = 6.000 \sin 75^\circ/\sin 50^\circ
\approx \qty{7.566}{km}$. In $ACD$: the angle at $D$ is $180 - 50 - 68 =
\qty{62}{\degree}$, and $AD = 7.566 \sin 68^\circ/\sin 62^\circ \approx
\qty{7.945}{km}$.
\end{solution}

\begin{solution}{exo:g10:measuring-the-earth:14}
$\cos\varphi = 1/2$ gives $\varphi = \qty{60}{\degree}$; $\cos\varphi = 0.1$
gives $\varphi \approx \qty{84.3}{\degree}$. A degree of longitude at
latitude $\varphi$ is the degree on the equator multiplied by $\cos\varphi$:
the ratio depends only on the angle, not on the radius.
\end{solution}

\begin{solution}{exo:g10:measuring-the-earth:15}
At the equator: $111.195 - 110.574 = \qty{0.621}{km} = \qty{621}{m}$; at the
pole: $111.694 - 111.195 = \qty{0.499}{km} = \qty{499}{m}$. Errors of several
hundred metres per degree, adding up degree after degree, are far larger than
the metre: positions converted from latitudes and longitudes must use the
ellipsoid.
\end{solution}

\begin{solution}{pb:g10:measuring-the-earth:1}
\textbf{1.} $31.20 - 24.09 = \qty{7.11}{\degree}$.

\textbf{2.} $7.11/360 \approx 0.0198 = 1/50.6$, very close to
Eratosthenes' $1/50$.

\textbf{3.} $C = 788 \times 360/7.11 \approx \qty{39900}{km}$.

\textbf{4.} The meridian arc between the parallels: the difference of the
noon angles is the difference of latitude, which measures only the
north--south separation of the two places; the east--west part of the
distance plays no role.

\textbf{5.} At local noon the Sun is in the plane of the meridian, the same
for both places; on the same day it has the same position among the stars,
so that the two angles differ only because of the curvature of the Earth.

\textbf{6.} $180 - 58 - 66 = \qty{56}{\degree}$.

\textbf{7.} $AC = 12.000 \sin 66^\circ/\sin 56^\circ = 12.000 \times
0.9135/0.8290 \approx \qty{13.223}{km}$.

\textbf{8.} $BC = 12.000 \sin 58^\circ/\sin 56^\circ \approx \qty{12.275}{km}$.

\textbf{9.} The angle at $D$ is $180 - 47 - 75 = \qty{58}{\degree}$;
$AD = 13.223 \sin 75^\circ/\sin 58^\circ \approx \qty{15.061}{km}$.

\textbf{10.} $15.061 \cos 8^\circ \approx \qty{14.914}{km}$.

\textbf{11.} $14.914/111.13 \approx \qty{0.1342}{\degree}$, that is
$0.1342 \times 60 \approx 8.05$ minutes of arc.

\textbf{12.} $111.51 - 110.57 = \qty{0.94}{km}$: \qty{940}{m}.

\textbf{13.} The vertical turns more slowly with distance near the pole: the
surface is less curved there. The Earth is flattened at the poles and bulges
at the equator.

\textbf{14.} $a - b = \qty{21.385}{km}$; $f = 21.385/6378.137 \approx 1/298$.

\textbf{15.} The plumb line is perpendicular to the surface of the
ellipsoid. On an ellipse the perpendicular to the curve passes through the
centre only at the equator and at the poles; at \qty{45}{\degree} it crosses
the axis below the centre.

\textbf{16.} $10\,000\,000$ metres, that is \qty{10000}{km}.

\textbf{17.} $10\,001\,966/10\,000\,000 \approx \qty{1.0001966}{m}$.

\textbf{18.} Too short: the true ten-millionth is $1.0001966$ of the 1799
metre, so the 1799 metre is short by about $1.97 \times 10^{-4}$ of its
length, 0.02\,\%.

\textbf{19.} Too short. The bar is one ten-millionth of the quarter meridian
as computed in the 1790s; since the bar is shorter than one ten-millionth of
the true quarter meridian, the computed quarter meridian was shorter than the
true one.

\textbf{20.} $1.0001966 - 1 = 0.0001966$ m: the first metre was about
\textbf{\qty{0.2}{mm} too short} --- about two sheets of paper over a metre,
an astonishing precision for a survey made with telescopes and angles across
a whole country.
\end{solution}
